\chapter{Partie de Rayan}
\label{chap:partie-rayan}

\section{Première section}
Texte de la section. Référence bibliographique : \cite{exemple}.

Exemple de formule :
\begin{equation}
    E = mc^2
    \label{eq:einstein}
\end{equation}

Exemple de tableau :
\begin{table}[H]
    \centering
    \begin{tabular}{lcc}
        \toprule
        Colonne A & Colonne B & Colonne C \\
        \midrule
        Valeur 1  & 10        & 20 \\
        Valeur 2  & 30        & 40 \\
        \bottomrule
    \end{tabular}
    \caption{Exemple de tableau}
    \label{tab:exemple}
\end{table}
